\documentclass[11pt,letterpaper]{article}
\usepackage[top=0.55in, bottom=0.4in, left=0.7in, right=0.7in]{geometry}
\usepackage{lmodern}
\usepackage{hyperref}
\hypersetup{colorlinks=false, urlcolor=black, pdfborder=0 0 0}
\setlength{\parindent}{0pt}
\setlength{\parskip}{0.65em}
\raggedright
\pagestyle{empty}

\begin{document}

\noindent
{\bfseries Ke-Cheng Hsu (Clio)}\\
Ph.D.\ Candidate, Department of Economics, HKUST\\
(852) 6840\,9753 \quad $|$ \quad \href{mailto:khsu@connect.ust.hk}{khsu@connect.ust.hk}\\
\href{https://cliohsu1997.github.io/}{cliohsu1997.github.io} \quad $|$ \quad \href{https://github.com/cliohsu1997}{github.com/cliohsu1997}

\vspace{0.4em}
\noindent
16 November 2026

\vspace{0.3em}
\noindent
Hiring Team\\
Canfield Securities Company Limited\\
Re: Junior Quantitative Research Analyst

\vspace{0.15em}
Dear Hiring Team,

I am writing to apply for the Junior Quantitative Research Analyst position. I am a Ph.D.\ candidate in Economics at HKUST (expected June 2027; available as early as February 2027), a Hong Kong permanent resident with the right to work in Hong Kong, a native Mandarin speaker, and professionally fluent in English. The role asks for quantitative research support on systematic strategies: collect and validate financial and market data, develop and test mathematical and statistical models, build Python research tools, use AI to speed implementation, and write clear, reproducible documentation. That is the skill set I have built through doctoral research.

Academic excellence and a strong quantitative foundation are the base. I completed a B.Sc.\ in Mathematics and Economics with a minor in Actuarial Mathematics at HKUST (CGA 3.9/4.3; Dean's List; Academic Achievement Medal, top 1\%), an M.Phil.\ in Economics (CGA 4.1/4.3), and I am finishing a Ph.D.\ in empirical industrial organization and digital economics. My coursework and research emphasize probability, statistics, structural modelling, GMM, Monte Carlo simulation, and computational implementation---the same mathematical and statistical toolkit used in quantitative investment research.

Programming, data infrastructure, and AI-assisted development are how I deliver. I build and maintain tested Python / R / C++ pipelines on AWS and HPC/SLURM for large proprietary panels: cleaning, linking, validating, and version-controlling data so research results are reproducible. I use AI tools every day for coding, automation, and research workflows, and I verify outputs with tests before relying on them. I learn new technical methods quickly and document models and results so others can follow the process.

My projects map to modelling, testing, and reporting. In my sole-authored job market paper on Taobao apparel return policy, I assemble a multi-source transaction dataset, estimate a structural model, run Monte Carlo and bootstrap checks, and compare counterfactual policies with clear performance metrics (a moderate restocking-fee menu raises social welfare by about 2 million RMB, roughly a 5\% gain, and seller profit by about 1.2 million RMB relative to status-quo policies). With Professors Kohei Kawaguchi and Kei Kawai, I estimate pricing and inventory under demand uncertainty on bookstore point-of-sale data spanning tens of billions of records, validate estimators with simulation, and deliver weekly analytical summaries on a fixed deadline. Earlier interbank-crisis modelling in MATLAB gave me exposure to financial-network dynamics, positions, and risk-style cascade logic. Across these projects I write technical documentation and short research notes that keep analysis consistent and auditable.

Clear communication closes the loop. As a lecturer and teaching assistant at HKUST (4.53/5 and 4.72/5), and as Teaching Assistant Coordinator in the Business School, I explain quantitative results to non-specialists and keep written outputs organized. I would welcome the chance to support systematic macro research at Canfield---data work, model development and backtesting-style evaluation, Python tooling, AI-assisted implementation, and analytical reporting. My resume is attached. Thank you for your consideration.

\vspace{0.4em}
Sincerely,\\[0.35em]
Ke-Cheng Hsu (Clio)

\end{document}
